% Separate post-CP17 source: not part of the frozen checkpoint-17 manuscript.
\begin{corollary}[Proved: surface-detected Brauer-zero growth on complete local punctures]
\label{D:local-Brauer-zero-surface-transfer}
Let $(R,\mathfrak m,k)$ be a complete excellent normal local domain
of dimension $d\ge2$, of equal characteristic $p>0$, with an
algebraically closed coefficient field $k$. Put
$U=\operatorname{Spec}R\setminus\{\mathfrak m\}$.
There are finitely many fixed complete normal local surface
$k$-algebras $R_1,\ldots,R_t$, with residue field $k$, and
morphisms
\[
 h_i:U_i=\operatorname{Spec}R_i\setminus\{\mathfrak m_i\}
                         \longrightarrow U
\]
with the following property. There is a constant $C_R$ such that,
for every $q=p^f$ and $r\ge1$,
\[
 \#\left(\left\{\rho:\pi_1(U)\longrightarrow
       \operatorname{PGL}_r(\mathbf F_q):
       \operatorname{br}(h_i^*\rho)=0\text{ in }
              H^2(U_i,\mathbf G_m)\text{ for every }i
        \right\}/\operatorname{PGL}_r(\mathbf F_q)\right)
                       \le q^{C_Rr^2}.                 \tag{D.LBT.1}
\]
In particular the bound holds for all representations with
$\operatorname{br}(\rho)=0$ in $H^2(U,\mathbf G_m)$.
The scalar centre is not required to have bounded order.

Fix $\epsilon>0$ and, for each isomorphism type $H$ in a
collection of nonabelian finite simple groups, choose one
embedding $H\hookrightarrow\operatorname{PGL}_{r_H}(\mathbf F_{q_H})$
with $q_H$ a power of $p$ and
$|H|\ge q_H^{\epsilon r_H^2}$. The quotient of $\pi_1(U)$
formed by all distinct simple quotient kernels whose chosen
projective representations satisfy the displayed surface
Brauer-vanishing condition is topologically finitely generated.
\end{corollary}

\begin{proof}
We use the fixed slice in
Proposition~\ref{D:local-surface-reduction}, and recall its
geometric maps in order to keep the Brauer condition precise.
Choose $d-2$ members of a system of parameters, with ideal $I$,
and set
\[
       C=(R/I)_{\mathrm{red}},\qquad
       W=\operatorname{Spec}C\setminus\{\mathfrak m_C\}.
\]
For $d=2$, take $I=0$. Every irreducible component of $C$
has dimension two. Its finite normalization is a disjoint
union of complete normal local surface domains $R_i$.
Their residue fields are finite over $k$ and therefore equal
to $k$; the original coefficient embedding supplies their
coefficient fields. If
\[
            W'=\coprod_{i=1}^t U_i,
\]
then $W'\to W$ is finite and surjective. The maps $h_i$ are
the composites $U_i\to W\to U$. All these choices are fixed
independently of $q$, $r$, and the representation.

The homomorphism $\pi_1(W)\to\pi_1(U)$ is surjective.
Indeed a connected finite etale cover of $U$ normalizes to
a finite normal domain $B$ over $R$. Completeness and
henselianity make $B$ local and complete, with dimension $d$.
The ideal $IB$ has at most $d-2$ generators. The connectedness
and cohomological-dimension lemmas
\citep[Lemmas 52.14.1 and 51.4.3, Tags 0ECR and 0DX9]{C_StacksLocalSlice}
show that $\operatorname{Spec}(B/IB)$ minus its closed point
is connected. Here the connectedness dimension of the domain
$B$ is $d$, and $d-2<d-1$ is the required strict inequality.
Removing nilpotents does not affect connectedness or finite
etale covers. Thus every original connected cover remains
connected after pullback to $W$, which proves the surjection.
In particular, two representations of $\pi_1(U)$ which become
projectively conjugate on $W$ were already projectively
conjugate on $U$.

Put $G_{r,q}=\operatorname{PGL}_r(\mathbf F_q)$ and let
\[
                     a=\#\pi_0(W'\times_W W').
\]
This is a finite integer depending only on the fixed slice
and normalization. For each $i$, Theorem~\ref{D:local-Brauer-zero-projective-growth}
bounds the relevant $G_{r,q}$-torsor classes on $U_i$ by
$q^{C_i r^2}$, where $C_i$ depends only on $R_i$.
Consequently their tuples on $W'$ have at most
$q^{(\sum_i C_i)r^2}$ possibilities.

We give the exact bound for descent fibres. Fix one such
tuple, regarded as a $G_{r,q}$-torsor $P'$ on $W'$. A torsor
on $W$ whose pullback has this isomorphism class can be
identified with $P'$ upstairs and hence supplies a descent
isomorphism between the two pullbacks of $P'$ to
$W'\times_W W'$. On each connected component, a morphism
between finite etale schemes is determined by its value at
one geometric point. If these are equivariant isomorphisms
of $G_{r,q}$-torsors, there are at most $|G_{r,q}|$ possible
values. Thus there are at most $|G_{r,q}|^a$ descent
isomorphisms, before imposing the cocycle equations or
identifying isomorphic descent data. Effective finite-proper
descent is \citet[Expose IX, Theorem 4.7 and its proof]{C_SGA1}.
It follows that each pullback fibre on torsor isomorphism
classes has cardinality at most $|G_{r,q}|^a$.
This bound uses ambient projective torsors and therefore
introduces no comparison between simple-group conjugacy
and ambient projective conjugacy.

Since $|G_{r,q}|\le q^{r^2}$, restriction to $W$ and the
preceding estimates give
\[
        \#\{\text{classes in (D.LBT.1)}\}
              \le q^{(a+\sum_i C_i)r^2}.
\]
Brauer obstruction is functorial under each actual morphism
$h_i$, so original Brauer-zero representations are among
the counted classes. This proves the quantitative assertion.

For the final assertion, distinct simple quotient kernels
give distinct ambient projective conjugacy classes after
the chosen embedding. Their multiplicity for type $H$ is
therefore at most $|H|^{C_R/\epsilon}$. The generating-tuple
estimate and Goursat/compactness argument in the proof of
Corollary~\ref{D:Brauer-zero-selected-simple-product} apply
verbatim, with this fixed constant, and yield a uniform
number of generators for the resulting profinite product.
\end{proof}

\begin{remark}[Scope of the transfer]
Only pullback of Brauer classes is used. No injection from
$H^2(\pi_1(U),-)$ into the cohomology of the detecting
surfaces, and no injectivity of the Brauer pullback, is
asserted. The conclusion concerns complete equicharacteristic
normal local punctures; it gives no mixed-characteristic,
arbitrary-open, or full fundamental-group finite-generation
statement. It is a closure of the established Brauer-zero
reduction, not a separate advance on the remaining genuinely
nonzero-Brauer or general crown problems.
\end{remark}
